\documentclass[11pt,letterpaper]{article}
\usepackage[T1]{fontenc}
\usepackage{lmodern,amsmath,amssymb,amsthm,mathtools}
\usepackage[margin=1in]{geometry}
\usepackage{microtype,enumitem,longtable,booktabs}
\usepackage[hyphens]{url}
\usepackage[colorlinks=true,linkcolor=black,citecolor=black,urlcolor=blue]{hyperref}
\hypersetup{pdftitle={Irreducibility of a family of shifted binomial compositions},pdfauthor={AI-assisted mathematical research},pdfsubject={Scoped irreducibility theorem and formal verification; originality unestablished}}
\setlength{\parskip}{0.25em}
\setlength{\parindent}{1.2em}
\setlist{nosep}
\newtheorem{theorem}{Theorem}[section]
\newtheorem{lemma}[theorem]{Lemma}
\newtheorem{proposition}[theorem]{Proposition}
\newtheorem{corollary}[theorem]{Corollary}
\theoremstyle{remark}\newtheorem{remark}[theorem]{Remark}
\newcommand{\Q}{\mathbb Q}
\newcommand{\Z}{\mathbb Z}
\newcommand{\R}{\mathbb R}
\newcommand{\C}{\mathbb C}
\newcommand{\F}{\mathbb F}
\newcommand{\Tr}{\operatorname{Tr}}
\newcommand{\Res}{\operatorname{Res}}
\newcommand{\N}{\operatorname{N}}
\newcommand{\OO}{\mathcal O}
\newcommand{\lean}[1]{{\small\urlstyle{tt}\nolinkurl{#1}}}
\title{Irreducibility of a family of shifted binomial compositions}
\author{Mathematical research draft prepared with AI assistance\\[0.3em]\small Project series: Noah's GPT-6 Astra math problems}
\date{12 October 2026}
\begin{document}
\maketitle
\begin{abstract}
Let $5\le p<q$ be primes and let $\kappa\in\Z$ satisfy $|\kappa|>1$ and
$v_\ell(\kappa)<p$ at every prime divisor $\ell$ of $\kappa$.
We give an argument that $(X^q-A)^p+\kappa$ is irreducible over $\Q$ for every
$A\in\Z$. The argument first descends a hypothetical affine power root from
the full ring of integers to the pure power order. A small-height
Fourier and congruence argument, and a large-height second Fourier mode
estimate with an integral resultant, then exclude that root.
The statement permits primes $p$ or $q$ dividing $\kappa$ and permits all
signs of $A$ and $\kappa$.
An accompanying Lean development and its verification boundary are described
explicitly. Originality has not been established. This result does not prove
the fixed-gap Pillai conjecture.
\end{abstract}

\section{Statement and scope}

For a nonzero integer $a$ and a rational prime $\ell$, $v_\ell(a)$ denotes
the exponent of $\ell$ in $|a|$. All irreducibility statements in this paper
refer to polynomials over $\Q$, unless another field is specified.

\begin{theorem}\label{thm:main}
Let $p,q$ be primes with $5\le p<q$. Let $\kappa,A\in\Z$ satisfy
\begin{equation}\label{eq:hyp}
 |\kappa|>1,\qquad
 v_\ell(\kappa)<p\quad\text{for every prime }\ell\mid\kappa.
\end{equation}
Then
\[
                 (X^q-A)^p+\kappa
\]
is irreducible over $\Q$.
\end{theorem}

In particular, the parameter is $p$th-power-free in the literal integer
sense. No assumption $p\nmid\kappa$ or $q\nmid\kappa$ occurs. The constant
$A$ is an arbitrary integer, including zero. No perfect-power condition on
$A^p-\kappa$ is assumed.

The proof has two main parts. Section~\ref{sec:order} proves a rigidity
theorem for integer coefficients in an irreducible pure power order. Section
\ref{sec:local} proves that the coefficients of a hypothetical root in the
pure field are indeed integers under \eqref{eq:hyp}. Section
\ref{sec:assembly} supplies the signs and the composition argument.
All orders used below are distinguished from full rings of integers.

\paragraph{Research status.}
The argument is a scoped auxiliary result arising from an investigation of
Pillai's conjecture. Its originality remains unestablished after the bounded
literature comparison in Section~\ref{sec:prior}. Earlier exploratory
arguments involving discriminant indices, phase clusters and successive
$q$-adic congruences are unnecessary for this proof. They are not used as
unproved bridges. The formal-verification status in Section~\ref{sec:verification}
must be read separately from both originality and the full Pillai problem.

\section{Rigidity in an integral pure power order}\label{sec:order}

\begin{theorem}\label{thm:order}
Let $n\ge5$ be an integer, let $k>1$ be an integer, and suppose that
$T^n-k$ is irreducible over $\Q$. Let $\rho>0$ satisfy $\rho^n=k$.
If $q>n$ is prime, $A>0$ is an integer and $s\in\{1,-1\}$, then there is no
$\gamma\in\Z[\rho]$ satisfying
\begin{equation}\label{eq:orderpower}
                         \gamma^q=A-s\rho.
\end{equation}
\end{theorem}

Here $n$ need not be prime, and $q$ is allowed to divide $k$. Membership in
$\Z[\rho]$, rather than only integrality in $\Q(\rho)$, is a hypothesis.
We prove the theorem by two exhaustive height cases.

Write
\[
 G(T)=\sum_{j=0}^{n-1}c_jT^j\in\Z[T],\qquad \gamma=G(\rho),
\]
and let $\zeta=\exp(2\pi i/n)$. Irreducibility and monic polynomial division
show that $T^n-k$ divides $G(T)^q-A+sT$ in $\Z[T]$. Thus, with indices
understood modulo $n$,
\begin{equation}\label{eq:conjugates}
 \gamma_i=G(\rho\zeta^i),\qquad
 \gamma_i^q=A-s\rho\zeta^i,\qquad
 |\gamma_i|\le R:=(A+\rho)^{1/q}.
\end{equation}
The use of all these conjugate equations is justified by divisibility,
not assumed independently.

\subsection{The small-height case}

First suppose $q\mid k$. Reducing the divisibility modulo $q$ and then
modulo $T^n$, Frobenius gives
\[
 G(T)^q\equiv c_0\pmod{q,T^n},
\]
because $jq\ge q>n$ for $j\ge1$. This contradicts the nonzero coefficient
$-s$ of $T$ in $A-sT$. Consequently we may assume $q\nmid k$.

Suppose now that
\begin{equation}\label{eq:small}
                            R<q\rho.
\end{equation}
Finite Fourier inversion in \eqref{eq:conjugates} yields
\[
 c_j\rho^j=\frac1n\sum_{i=0}^{n-1}\gamma_i\zeta^{-ij},\qquad
 |c_j|\le R/\rho^j.
\]
Since $\rho>1$, every $j\ge1$ has $|c_j|<q$.
In $\F_q[T]/(T^n-k)$ the same power equation becomes
\begin{equation}\label{eq:frob}
 \sum_{j=0}^{n-1}\overline{c_j}\,
 \overline{k}^{\lfloor jq/n\rfloor}T^{jq\bmod n}=\overline A-\overline sT.
\end{equation}
Here $\gcd(q,n)=1$ and $\overline k\ne0$. The residues $jq\bmod n$ are
therefore a permutation, and all nonconstant coefficients except possibly
$c_{j_0}$ vanish modulo $q$, where $j_0q\equiv1\pmod n$ and $1\le j_0<n$.
Their absolute-value bounds make them zero as integers. It follows that
\[
                  G(T)=a+bT^{j_0},\qquad \gcd(j_0,n)=1.
\]
The following elementary geometric lemma excludes this possibility.

\begin{lemma}\label{lem:circle}
Let $n\ge5$, $\gcd(j,n)=1$, $\rho,A>0$, $s\in\{1,-1\}$,
$a,b\in\R$, and let $m>1$ be an integer. The simultaneous identities
\[
 |a+b\rho^j\zeta^{ij}|^{2m}=|A-s\rho\zeta^i|^2
 \quad(0\le i<n)
\]
are impossible.
\end{lemma}
\begin{proof}
Put $S=\{\cos(2\pi i/n):0\le i<n\}$. It has maximum $1$ and at least
three distinct elements. The permutation $i\mapsto ij$ of $\Z/n\Z$ makes the cosine values
$\cos(2\pi ij/n)$ range over $S$, with the value at $i=0$ still equal to $1$.
Define on $[-1,1]$
\[
 f(x)=(a^2+b^2\rho^{2j}+2ab\rho^j x)^m,
 \qquad h(x)=A^2+\rho^2-2sA\rho x.
\]
The asserted equations imply $f(S)=h(S)$ and $f(1)=h(1)$.
If $ab=0$, the first function is constant and the second is not, which is
impossible. Otherwise the base of $f$ is nonnegative throughout $[-1,1]$,
as its minimum is $(|a|-|b|\rho^j)^2$. Hence $f$ is strictly monotone and
strictly convex, including when that minimum is zero at an endpoint.
The affine function $h$ is strictly monotone.

Their monotonic orientations must agree: if opposite, their common value
at the maximum argument $1$ would be both the minimum and maximum of the
same nontrivial image set. Sorting the two equal image sets then gives
$f(x)=h(x)$ at every $x\in S$. At three ordered elements of $S$ this
contradicts strict convexity of $f$ and affinity of $h$.
\end{proof}

Applying Lemma~\ref{lem:circle} to \eqref{eq:conjugates} proves the
small-height exclusion. No assertion $j_0=1$ or order-discriminant equality
is required.

\subsection{An integral second Fourier coefficient}

For the complementary height range we use a slightly more general integer
$m>n$. Assume \eqref{eq:conjugates} with $q$ replaced by $m$, where
$A>0$ may now be rational. Define the integer polynomial
\begin{equation}\label{eq:autocorr}
 H(T)=\sum_{\substack{0\le j,t<n\\j-t\equiv2\ (\mathrm{mod}\ n)}}
                      c_jc_tT^{j+t},\qquad
 \beta_h=H(\rho\zeta^h).
\end{equation}
Finite orthogonality gives the exact identities
\begin{align}
 \beta_h&=\frac1n\sum_{i=0}^{n-1}
                    \gamma_{h+i}\gamma_{h-i}\zeta^{-2i},\label{eq:twist}\\
 \beta_0&=\frac1n\sum_{i=0}^{n-1}|\gamma_i|^2\zeta^{-2i}.
                                                        \label{eq:mode}
\end{align}
The second uses real coefficients, so $\gamma_{-i}=\overline{\gamma_i}$.
The first and $|\gamma_i|\le R$ imply
\begin{equation}\label{eq:allbounds}
                              |\beta_h|\le R^2
                  \quad\text{for every }h.
\end{equation}
In particular, the small Fourier quantity is an algebraic integer in a
field of degree $n$, not in a growing field of independently chosen radicals.

\begin{lemma}\label{lem:taylor}
Let $n\ge5$, $m>n$, $a=1/m$, $s\in\{1,-1\}$ and $0<u<1/1000$.
Put $x_i=\cos(2\pi i/n)$. Then
\[
 0<\left|\frac1n\sum_{i=0}^{n-1}
             (1-2sx_i u+u^2)^a\zeta^{-2i}\right|<a u^2.
\]
\end{lemma}
\begin{proof}
For $x\in[-1,1]$, write $h(v)=1-2sxv+v^2$ and $g(v)=h(v)^a$.
On $0\le v\le1/4$,
\[
 h(v)\ge(1-v)^2\ge9/16>1/2,
 \quad |h'(v)|\le5/2<3,\quad h''(v)=2.
\]
Since $0<a\le1/6$, we have $h^{a-3}\le8$, $h^{a-2}\le4$,
$|(a-1)(a-2)|\le2$ and $|a-1|\le1$. The identity
\[
 g'''=a(a-1)(a-2)h^{a-3}(h')^3
                         +3a(a-1)h^{a-2}h'h''
\]
therefore gives $|g'''|\le504a$. Taylor's theorem implies, uniformly in $x$,
\[
 g(u)=1-2asxu+[a+2a(a-1)x^2]u^2+e_x,
                      \qquad |e_x|\le84a u^3.
\]
Explicitly, $x_i=(\zeta^i+\zeta^{-i})/2$ and
$x_i^2=(\zeta^{2i}+2+\zeta^{-2i})/4$. Multiplication by $\zeta^{-2i}$
and summation leave only the constant term $1/4$ in the last expression:
none of $1,2,3,4$ is a nonzero multiple of $n\ge5$. The mode-two
coefficients of $1,x_i,x_i^2$ are therefore $0,0,1/4$. Thus the expression in the lemma
is
\[
               \frac{a(a-1)}2u^2+E,\qquad |E|\le84a u^3<a u^2/6.
\]
The leading term is negative and has absolute value between $a u^2/3$
and $a u^2/2$. These bounds prove both strict inequalities.
\end{proof}

\subsection{The large-height case}

Suppose
\begin{equation}\label{eq:large}
                        A+\rho\ge(m\rho)^m,
                  \qquad m>n\ge5.
\end{equation}
Set $R=(A+\rho)^{1/m}$ and $u=\rho/A$.
Since $\rho>1$ and $m\ge6$,
\[
                 0<u<\frac1{m^m-1}<\frac1{1000}.
\]
The actual absolute-value equations, without selecting any complex root
branches, give
\[
 |\gamma_i|^2=A^{2/m}(1-2sx_i u+u^2)^{1/m}.
\]
Lemma~\ref{lem:taylor} and \eqref{eq:mode} consequently imply
\begin{equation}\label{eq:onebound}
                0<|\beta_0|<\frac{\rho^2 R^2}{mA^2}.
\end{equation}
Because $T^n-k$ is irreducible and $H(\rho)=\beta_0\ne0$, the two integer
polynomials $T^n-k$ and $H$ have a nonzero integer resultant. The first is
monic with distinct roots $\rho\zeta^i$, so
\begin{equation}\label{eq:resultant}
           0\ne\Res(T^n-k,H)=\prod_{i=0}^{n-1}\beta_i\in\Z.
\end{equation}
On the other hand, $R\ge m\rho>1$ and
$R^m\ge(m\rho)^m>2\rho$, whence $A=R^m-\rho>R^m/2$.
Using \eqref{eq:allbounds}, \eqref{eq:onebound} and $n\le m-1$ gives
\begin{align*}
 \left|\prod_i\beta_i\right|
 &<\frac{\rho^2R^{2n}}{mA^2}
 <\frac{4\rho^2R^{2n}}{mR^{2m}}
 \le\frac{4\rho^2}{mR^2}
 \le\frac4{m^3}<1.
\end{align*}
This contradicts \eqref{eq:resultant}. Taking $m=q$ covers exactly the
complement of \eqref{eq:small}, and proves Theorem~\ref{thm:order}.

\section{Descent from the pure field to the power order}\label{sec:local}

\begin{proposition}\label{prop:descent}
Let $p,q$ be primes with $3\le p<q$, and let $\kappa,A\in\Z$ satisfy
\eqref{eq:hyp}. Let $\theta^p=\kappa$ and $K=\Q(\theta)$.
If $\gamma\in K$ satisfies $\gamma^q=A-\theta$, then
\[
                   A\ne0,\qquad \gamma\in\Z[\theta].
\]
\end{proposition}

We prove all the local inputs and then intersect rational coefficient
conditions. The proof applies to signed $\kappa$ and $A$.

\subsection{Normalization and degree}

Choose a prime $\ell\mid\kappa$, and write
\[
       \kappa=\ell^r U,\qquad 0<r<p,\qquad U\in\Z,\quad\ell\nmid U.
\]
These are conclusions from the exact valuation, not additional hypotheses.
Choose integers $1\le a<p$, $b\ge0$ such that $ar=bp+1$. In any
$\ell$-adic field containing a root $t^p=\kappa$, put $\pi=t^a/\ell^b$.
Then
\begin{equation}\label{eq:normalization}
             \pi^p=\ell U^a,\qquad \pi^r=tU^b,
                         \qquad t=U^{-b}\pi^r.
\end{equation}
The polynomial $T^p-\ell U^a$ is Eisenstein over $D=\Z_\ell$.
It follows that $\Q_\ell(t)=\Q_\ell(\pi)$ has degree $p$, so
$T^p-\kappa$ is irreducible locally and hence over $\Q$.
The full integral closure of $D$ in this field is
\begin{equation}\label{eq:closure}
                              S=D[\pi].
\end{equation}
Here is the local valuation argument for this standard Eisenstein fact.
Normalize the extended $\ell$-adic valuation by $v(\ell)=1$.
For a root $\pi$ of any degree-$d$ monic Eisenstein polynomial over
$\Z_\ell$, its defining equation forces $v(\pi)=1/d$.
Indeed integrality gives $v(\pi)\ge0$; value zero would make the leading
term uniquely smallest, while a positive value different from $1/d$ would
make either the leading or constant term uniquely smallest. Both are
impossible in a vanishing sum.
For $z=\sum_{j=0}^{d-1}b_j\pi^j$ with $b_j\in\Q_\ell$, its nonzero
summands have valuations $v(b_j)+j/d$, distinct modulo $\Z$. Thus
$v(z)=\min_j(v(b_j)+j/d)$. An integral $z$ has $v(z)\ge0$, which forces
$v(b_j)\ge0$ for every $j$, because these coefficient valuations are
integers and $j<d$. Hence $z\in\Z_\ell[\pi]$. The reverse inclusion
follows from integrality of $\pi$. This proves \eqref{eq:closure} and
also justifies the other Eisenstein closure applications below.
No tameness assumption occurs, including when $\ell=p$.

The element $\gamma$ is an algebraic integer: it satisfies the monic
integer polynomial $(X^q-A)^p+\kappa$. Write its unique reduced expansion
\begin{equation}\label{eq:coords}
                   \gamma=\sum_{j=0}^{p-1}c_j\theta^j,
                              \qquad c_j\in\Q.
\end{equation}
Set $G(T)=\sum_{j=0}^{p-1}c_jT^j\in\Q[T]$.
Every equation in this expansion propagates to any root of $T^p-\kappa$
over a characteristic-zero field, by divisibility by that irreducible
polynomial. At the local root $t$, its value $x=\sum c_jt^j$ is integral
over $D$, because $x^q=A-t$ and $t$ is integral. Hence $x\in S$ by
\eqref{eq:closure}.

\subsection{Primes dividing the parameter}

Retain $\ell,r,U,D,S$ above. Write $c=U^{-b}\in D^\times$, so
$t=c\pi^r$. The actual reduction of the Eisenstein power basis is
\[
            S/\ell S\simeq\F_\ell[T]/(T^p),\qquad\pi\longmapsto T.
\]
If $\ell\mid A$, the constant coefficient of the reduction of
$x^q=A-c\pi^r$ forces the constant coefficient of $\overline x$ to vanish.
Thus $\overline x\in(T)$ and $\overline x^q=0$, since $q>p$.
But $-\overline cT^r\ne0$ for $0<r<p$, a contradiction. We have proved
\begin{equation}\label{eq:unitA}
                              \ell\nmid A.
\end{equation}
In particular, any such parameter prime proves $A\ne0$.

If $\ell=q$, Frobenius gives an even stronger contradiction: every $q$th
power in $\F_q[T]/(T^p)$ is constant, whereas
$\overline A-\overline cT^r$ is not. Consequently a hypothetical affine
power equation automatically forces $q\nmid\kappa$.
This exclusion is why Theorem~\ref{thm:main} need not assume it.

Suppose $\ell\ne q$, and put $R=D[t]\subset S$, $J=\pi S$.
The map $S\to\F_\ell$ taking the constant coefficient modulo $\ell$ has
kernel $J$. Indeed $\ell=\pi^p U^{-a}$; a constant coefficient divisible
by $\ell$ is therefore divisible by $\pi$ in $S$, as are all other terms.
By \eqref{eq:unitA}, the residue of $x$ is nonzero. Choose
$a_0\in D^\times$ with that residue. Then $x-a_0\in J$ and
$q a_0^{q-1}\in D^\times\subset R^\times$.

A finite conductor inclusion suffices:
\begin{equation}\label{eq:conductor}
                    J^N\subset R,\qquad N=(p-1)(r-1).
\end{equation}
To see it explicitly, every integer $m\ge N$ is $ir+jp$ with $i,j\ge0$.
Choose $0\le i<p$ so that $ir\equiv m\pmod p$; then
$m-ir\ge N-(p-1)r=-(p-1)>-p$, so the multiple $m-ir$ of $p$ is nonnegative.
Equations \eqref{eq:normalization} imply $\pi^{ir+jp}\in R$.
Since $1,\pi,\ldots,\pi^{p-1}$ span $S$ over $D$, this proves
\eqref{eq:conductor}. When $r=1$, it says $S=R$ directly.

\begin{lemma}[Finite Newton descent]\label{lem:newton}
Let $R$ be a subring of a commutative ring $S$, let $J$ be an ideal of $S$,
and suppose $J^N\subset R$ for an integer $N\ge1$.
Let $f\in R[T]$, $f(x)=0$ for $x\in S$, and $x-a_0\in J$ with $a_0\in R$.
If $f'(a_0)$ is a unit of $R$, then $x\in R$.
\end{lemma}
\begin{proof}
Keep $d=f'(a_0)^{-1}\in R$ fixed and define
$a_{m+1}=a_m-df(a_m)$. All $a_m$ lie in $R$.
Inductively, $e_m=x-a_m\in J^{m+1}$.
Polynomial Taylor expansion gives
$f(a_m)=-f'(a_m)e_m+e_m^2z_m$ for some $z_m\in S$.
Also $f'(a_m)-f'(a_0)\in J$ because $a_m-a_0\in J$.
Hence
$e_{m+1}=(1-df'(a_m))e_m+d e_m^2z_m\in J^{m+2}$.
After $N-1$ iterations, $x-a_{N-1}\in J^N\subset R$, proving the claim.
This argument neither assumes completeness nor asserts the existence of a
new root.
\end{proof}

For $r=1$ we already have $x\in S=R$. For $r>1$, $N\ge1$; apply
the lemma to $f(T)=T^q-(A-t)$, with the data above, to obtain $x\in R$.
Since $t$ has degree $p$ over $\Q_\ell$, reduction modulo its monic
polynomial and uniqueness of the reduced expansion imply
\begin{equation}\label{eq:paramcoords}
                              c_j\in\Z_\ell\quad(0\le j<p).
\end{equation}
This includes the case $\ell=p$ when $p\mid\kappa$.

\subsection{The prime exponent when it is a parameter unit}

It remains to treat the prime $p$ if $p\nmid\kappa$.
We first justify the elementary local dichotomy used here. Choose an
integer $a\equiv\kappa\pmod p$, a $p$-adic unit. If
$p^2\nmid a^p-\kappa$, then the constant coefficient has valuation one,
so $(T+a)^p-\kappa$ is Eisenstein. Otherwise $p^2\mid a^p-\kappa$,
a case that includes $a^p-\kappa=0$. For odd $p$,
\[
 (a+pu)^p\equiv a^p+p^2 a^{p-1}u\pmod{p^3}.
\]
Choose $u\pmod p$ to make the right side congruent to $\kappa$ modulo $p^3$.
At $a+pu$, the polynomial $T^p-\kappa$ has value divisible by $p^3$ and
its derivative has valuation exactly one. Strong Hensel lifting therefore
produces $b\in\Z_p^\times$ with $b^p=\kappa$.
These alternatives exhaust the cases, since $a^p\equiv\kappa\pmod p$.

In the Eisenstein alternative, let $\lambda$ be a root of
$(T+a)^p-\kappa$ over $\Q_p$. Its integral closure is
$\Z_p[\lambda]=\Z_p[a+\lambda]$. The propagated equation and integrality
put $G(a+\lambda)$ in this order. The degree-$p$ basis then gives
$c_j\in\Z_p$.

In the split alternative choose $b$ as above and let $\zeta_p$ be a
primitive $p$th root of unity. The shifted cyclotomic polynomial
$\Phi_p(T+1)$ is Eisenstein, so $\Q_p(\zeta_p)$ has degree $p-1$ and
integral closure $\Z_p[\zeta_p]$. Propagation and integrality give
\[
  x=G(b)\in\Z_p,\qquad y=G(b\zeta_p)\in\Z_p[\zeta_p],
  \qquad x^q=A-b,\quad y^q=A-b\zeta_p.
\]
In the common residue field $\F_p$, $\overline x^q=\overline y^q$.
Because $q>p$ is prime, $\gcd(q,p-1)=1$, and the $q$th-power map is injective
on all of $\F_p$, including zero. Thus $\overline x=\overline y$.

Write $y=Q(b\zeta_p)$ with $Q\in\Z_p[T]$, $\deg Q<p-1$; this is possible
because $b$ is a unit. Then $x-Q(b)=pd$ for some $d\in\Z_p$, and
\[
                         P(T)=Q(T)+d\Phi_p(T/b)
\]
is an integral polynomial of degree less than $p$ with $P(b)=x$ and
$P(b\zeta_p)=y$. The polynomial $G-P$ has degree less than $p$ and vanishes
at $b\zeta_p$. Its minimal polynomial is a nonzero scalar multiple of
$\Phi_p(T/b)$ of degree $p-1$, so $G-P=e\Phi_p(T/b)$ for some $e\in\Q_p$.
Evaluating at $b$ gives $ep=0$, hence $G=P$. Therefore every $c_j$ is
$p$-integral in this alternative too.

\subsection{Trace bounds and the rational intersection}

For $K=\Q(\theta)$ the pure power basis gives
\[
 \Tr_{K/\Q}(\theta^m)=
 \begin{cases}0,&p\nmid m,\\p\kappa^{m/p},&p\mid m.\end{cases}
\]
The element $\gamma$ and all $\gamma\theta^{p-j}$ are integral. In
\eqref{eq:coords}, it follows that
\[
      \Tr(\gamma)=pc_0\in\Z,
 \qquad \Tr(\gamma\theta^{p-j})=p\kappa c_j\in\Z\quad(1\le j<p).
\]
Thus the reduced denominator of $c_0$ divides $p$, and that of each other
$c_j$ divides $p\kappa$. Every possible denominator prime has already been
excluded: parameter primes by \eqref{eq:paramcoords}, and $p$ by that same
argument when $p\mid\kappa$, or by the dichotomy when $p\nmid\kappa$.
A rational number belonging to $\Z_\ell$ has reduced denominator prime to
$\ell$. Therefore all the denominators equal one and
$\gamma\in\Z[\theta]$. Together with \eqref{eq:unitA}, this proves
Proposition~\ref{prop:descent}.

\section{All signs and the composition argument}\label{sec:assembly}

Suppose that $\gamma\in\Q(\theta)$ satisfies $\gamma^q=A-\theta$ under
Theorem~\ref{thm:main}'s hypotheses. Proposition~\ref{prop:descent} gives
$A\ne0$ and an integer polynomial representing $\gamma$ in $\theta$.
Put $k=|\kappa|$, $\rho=k^{1/p}>0$ and $s=\operatorname{sign}(\kappa)$.
The binomial $T^p-k$ is irreducible: it is $T^p-\kappa$ when $\kappa>0$,
and an odd sign substitution of that polynomial when $\kappa<0$.
The equation propagates to $\theta=s\rho$; absorbing $s^j$ into each
integer coefficient gives an element $\delta\in\Z[\rho]$ with
\[
                             \delta^q=A-s\rho.
\]
If $A>0$, Theorem~\ref{thm:order} contradicts this. If $A<0$, oddness of
$q$ gives $(-\delta)^q=(-A)-(-s)\rho$, again contradicting that theorem.
Thus no affine root exists in the degree-$p$ pure field for any $A\in\Z$.

For completeness, apply Capelli's composition criterion to
\[
                 f(Y)=Y^p+\kappa,\qquad g(X)=X^q-A.
\]
The polynomial $f$ is irreducible by odd sign substitution from
$Y^p-\kappa$. Let $\alpha$ be a root, and put $\theta=-\alpha$.
Then $\theta^p=\kappa$ and $\Q(\alpha)=\Q(\theta)$.
Over this actual simple field,
\[
                       g(X)-\alpha=X^q-(A-\theta).
\]
For an odd prime $q$ in characteristic zero, the prime-binomial criterion
states that $X^q-a$ is irreducible if and only if $a$ is not a $q$th power
in the field. To recall its short proof for $a\ne0$, a proper monic factor
of degree $0<d<q$ has roots among $\beta\zeta_q^i$, where $\beta^q=a$.
The product $u$ of those roots belongs to the base field, and $u^q=a^d$.
Choose integers $e,h$ with $ed+hq=1$. Then $(u^e a^h)^q=a$, as required.
The converse follows from a linear factor, and $a=0$ is immediate.
The affine exclusion above applies to every element of the actual simple
field: its degree-$p$ power basis supplies a rational polynomial
representative of degree less than $p$. Hence $g(X)-\alpha$ is irreducible.

One can see the needed direction of Capelli directly by taking a root
$\beta$ of $g(X)-\alpha$. Since $\alpha=g(\beta)$, the tower law gives
$[\Q(\beta):\Q]=[\Q(\beta):\Q(\alpha)]p=qp$.
The monic polynomial $f(g(X))$ of degree $pq$ has root $\beta$, and so is
its minimal polynomial. This proves Theorem~\ref{thm:main}.

\section{Relation to the fixed gap problem}

The norm implication is one-way. A $q$th root of $A-\theta$ in
$\Q(\theta)$ would give a $q$th-power rational norm, but a $q$th-power norm
does not imply the element has a root in that field. For example,
\[
              3^5-2^7=115=5\cdot23,
       \qquad 5^5-3^7=938=2\cdot7\cdot67.
\]
Both parameters satisfy the theorem's valuation hypotheses, including the
respective cases $p\mid\kappa$ and $q\mid\kappa$. The theorem therefore
asserts irreducibility of the corresponding degree-$35$ compositions,
while these perfect-power gaps genuinely exist.

The full fixed-gap Pillai conjecture asks for finitely many distinct
perfect-power value pairs at each fixed positive gap. Its moving-exponent
tail may occupy a nonzero Kummer class: a $q$th root of $A-\theta$ then lies
in a genuine extension of $\Q(\theta)$. The proof here excludes the zero
class only. In particular, neither a finite-height search nor the absence
of an affine root inside the pure field proves fixed-gap finiteness.

The integer resultant in Section~\ref{sec:order} depends crucially on
all conjugates coming from one integer polynomial in the same pure root.
For a radical lying in an extension, this representation is absent.
Neither the present theorem nor its formal verification supplies a
uniform exclusion or height bound for that remaining branch.

\section{Comparison with prior literature}\label{sec:prior}

Capelli's composition and prime-binomial criteria are standard;
Guersenzvaig~\cite{Guersenzvaig} gives a modern primary account.
The DVR Eisenstein closure theorem, Hensel lifting, traces and resultants
are likewise standard tools. The finite conductor proof and the elementary
Fourier estimates above are supplied explicitly rather than inferred from
general assertions about orders.

Harrington and Jones~\cite{HJ} study monogenic binomial compositions of
this general shape. The complete article was not accessible in the present
review, so it remains an explicit priority for a knowledgeable referee.
The accessible preprint of Jakhar, Kalwaniya and Yadav~\cite{JKY} assumes
the composition irreducible in its main index characterization. The newer
Choudhary, Pisolkar and Yadav preprint~\cite{CPY} likewise states
irreducibility as an assumption for its monogeneity criterion. Those
assumptions do not by themselves supply Theorem~\ref{thm:main}.
Ga\'al's related article~\cite{Gaal} concerns a degree-six family according
to its published abstract; the entire PDF was not used in this comparison.

Banerjee~\cite{Banerjee} gives variations on Capelli's theorem. Only a
publisher first-page preview was readable in this review. Its necessary
condition for reducibility of $f(X^q)$ is a divisibility congruence modulo
$q^2$. That quoted condition alone does not cover the present theorem.
For instance $p=5,q=7,A=49$ and
\[
           \kappa=49^5-2^7=282475121=23\cdot2953\cdot4159
\]
satisfy the valuation hypotheses and have a perfect-power norm. For
$f(Y)=(Y-A)^p+\kappa$, the condition
$f(Y)\mid f(Y^q)\pmod{q^2}$ holds. Indeed, more generally, if $q^2\mid A$
and $\kappa=A^p-B^q$, then $f(Y)\equiv Y^p-B^q\pmod{q^2}$, and the
remainder of $f(Y^q)$ is $B^{q^2}-B^q\equiv0\pmod{q^2}$.
This compares the displayed necessary condition only, not every result in
the unreviewed article.

An older composition reference, Butler's 1957 paper~\cite{Butler}, also
remains a full-text access gap. The bounded search found no theorem verified
to cover exactly all the quantifiers here, but absence of such a search
hit is not evidence sufficient to establish originality. No priority,
publication or new-result claim is made. The proof and its formal statement
are suitable for further specialist review independently of that question.

\section{Formal theorem mapping and verification}\label{sec:verification}

The development uses Lean 4.32.2 and mathlib commit
\begin{center}\small\texttt{905b95818eb32af7874a58b427f50c1711a5e96c}.\end{center}
The checked final root is
\begin{center}\small
\lean{PillaiScoped.}\allowbreak
\lean{irreducible_composition_of_prime_valuation_bounds}.
\end{center}
Its statement is literal polynomial irreducibility over $\Q$ under only
prime $p,q$, $5\le p<q$, $\kappa,A\in\Z$, $1<|\kappa|$ and
\[
 \forall\ell\in\mathbb N,\quad
 \ell\text{ prime}\ \Longrightarrow\ \ell\mid\kappa\
 \Longrightarrow\ \operatorname{padicValInt}(\ell,\kappa)<p.
\]
Because $1<|\kappa|$, Lean's convention for the valuation at zero is
irrelevant. There is no assumed irreducibility, local-field degree,
coefficient integrality, affine-exclusion callback or perfect-power gap
in this final statement.

\subsection{Declaration map}

The following map identifies the main checked components. The final status
is documented by the versioned whole-closure verification record accompanying
this file, separately from compilation of individual helper lemmas.
\begingroup\raggedright
\paragraph{Integral power order theorem}\mbox{}\par\noindent
Module \lean{IntegralPowerOrder}.\par\noindent Root:
\lean{PillaiIntegralPowerOrder.integral_power_order_exclusion_fin}.
It includes actual conjugate propagation, Frobenius coefficient reduction,
finite-circle convexity, Taylor nonvanishing, the resultant and the product
bound in Theorem~\ref{thm:order}.
\paragraph{Trace denominators}\mbox{}\par\noindent
Module \lean{TraceDenominators}.\par\noindent Root:
\lean{PillaiTrace.irreducible_denominator_dvd}.
It constructs the simple field and trace calculation rather than assuming
that the ambient field is generated by the root.
\paragraph{The unit parameter prime}\mbox{}\par\noindent
Module \lean{PadicCoefficientLift}.\par\noindent Root:
\lean{PillaiLocalCoefficients.padic_coefficients_of_power_equation}.
Its actual Eisenstein and split cyclotomic cases supply Section~\ref{sec:local}'s
$p\nmid\kappa$ branch.
\paragraph{Parameter primes and the ramified exponent}\mbox{}\par\noindent
Module \lean{ConcreteParameterCoefficients}.\par\noindent Roots:
\lean{PillaiLocalCoefficients.}\allowbreak
\lean{parameter_prime_coefficients_of_normalized_power_equation}
and
\lean{PillaiLocalCoefficients.}\allowbreak
\lean{no_normalized_power_equation_at_exponent_parameter}.
These construct the actual auxiliary field, full integral closure,
conductor descent and coefficient recovery. The first allows $\ell=p$;
the second has no degree bound on its rational polynomial.
\paragraph{Global arithmetic}\mbox{}\par\noindent
Modules \lean{IntegerValuationNormalization},
\lean{RationalLocalIntersection}, \lean{GlobalCoefficientAssembly}.\par\noindent Root:
\lean{PillaiGlobalCoefficients.integer_polynomial_of_power_equation}.
It derives $A\ne0$ and a literal reduced polynomial over $\Z$ from the
natural valuation hypotheses.
\paragraph{Signs and final composition}\mbox{}\par\noindent
Modules \lean{SignedRealExclusion}, \lean{CapelliPureFieldBridge},
\lean{ScopedPillaiIrreducibility}. The generic Capelli helper has an
explicit affine-exclusion input; the final root proves and discharges
that input through the actual global arithmetic and real exclusion.
\par\endgroup

\clearpage\subsection{Verification status}

At 04:05 UTC on 12 October 2026 the final scoped theorem and all 41 final-assembly
declarations passed a fresh sequential clean build and recursive
standard-axiom audit, with zero warnings. The frozen snapshot is dated
04:07 UTC. Its source manifest hash and final source hash are recorded in
the companion reproducibility file. The final clean peak process-group
RSS was 3,283,768 KiB. An earlier internal M3200 memory exception was
preserved; the successful final union used the separately authorized
M3500 cap and produced no new dependency builds.
At 04:14 UTC the final production theorem type and referenced definitions
matched a separately authored, reviewed specification of the requested
formula. Its complete serialized dependency closure passed the official
Lean 4.34.1 kernel and independent nanoda 0.4.19 checking, with 45,192
declarations and all 4,523 imported-object hashes matched.
The production closure SHA256 is
\begin{center}\small\nolinkurl{e708c9c22e754d2af9b542fd151819c37fa5c00bb169f4acc4c934e573d10652}.\end{center}
The sealed result JSON SHA256 is
\begin{center}\small\nolinkurl{18e6daf439c640cf8a26767339d145df0960793e70da7afcc25c3cd4e372effe}.\end{center}
Both hashes were independently recomputed when this manuscript was updated.

Earlier component closures also passed the same two checkers. The final
closure above contains every dependency needed for the stated theorem.
Each V2 check compares 26 mandatory primitive definitions and every present
permitted axiom type against an independently exported official-version
baseline. The permitted standard axioms are \lean{propext},
\lean{Classical.choice} and \lean{Quot.sound}; no \lean{sorryAx},
custom axiom or unproved mathematical bridge is accepted.

This is serialized full-closure replay in a patched official kernel plus
an independent checker, with a separate requested-type and
referenced-definition comparison. It used a non-isolated validation
adapter; it is not sandboxed comparator certification or external peer review.
The comparison's unproved type-specification fixture supplied only the
requested type, never a production proof. Its placeholder proof and
\lean{sorryAx} are excluded from all production imports and the proof-source
package. The exported production axioms are exactly the three standard ones.
Human reading of the exported theorem type is a separate semantic check.
An independent AI-assisted paper/source review found no mathematical gap;
that is not represented as external expert peer review.
Formal verification does not certify originality or prove Pillai's conjecture.

\subsection{Reproducibility}

The accompanying source package retains all 51 local production modules
and the declaration-audit module, selected by actual imported-object
provenance rather than filename recency. It includes source and object
hashes, a closed dependency order, build scripts, pinned toolchain and
mathlib metadata, imported-object inventory, printed statements, recursive
axiom audits, clean-build logs and compact checker records. The separate
admitted type-comparison fixture is excluded from the proof sources.

On 12 October 2026, from 04:28:57 to 04:33:09 UTC, all 52 packaged local
sources were compiled sequentially into a new empty object tree. No old
local proof object was present on the dependency search path; only pinned
upstream caches were reused. Every newly generated local object was
byte-identical to its historically validated counterpart. All 4,523
imported-object hashes matched, all 41 declaration audits used only the
three standard axioms, and there were no warnings or changed source bytes.
The sampled peak process-group RSS was 3,310,132 KiB. The documented memory
profile uses M3200 for helpers, M3300 for \lean{ActualPadicSplit}, and
M3500 for the final union/audit, with independent RSS, available-memory and
disk guards. The first M3200-only attempt and its memory-limit failure
are retained. The clean-build result SHA256 is
\begin{center}\small\nolinkurl{a804f775a3404703cc2924af65d9c3ce5ef35d83d826684252df691b8789103f}.\end{center}

At 04:35 UTC a new export from those freshly built objects was compared
with the previously validated production closure. All 294,553,707 bytes
were identical, with the same SHA256 and all imported-object hashes
rechecked. Thus the packaged-source rebuild is tied to the existing
requested-type comparison, full patched-kernel replay and independent
nanoda result; no inference from source resemblance alone is needed.

Lean's package-path lookup selects the first root containing a package
directory and does not fall back module by module. Consequently a partial
\lean{Mathlib} override directory must be merged with the remaining pinned
cache artifacts, rather than merely listed before another cache root.
The package's sequential source builder rejects dependency roots exposing
old local proof objects. Its pinned fresh-network installation recipe is
provided but has not been tested in a separate environment; the successful
clean-local reproduction used the recorded upstream caches.

For stronger replay, export the final root and its complete closure with
the authenticated original-version exporter, compare the primitive
baseline and imported hashes, and run the recorded official 4.34.1 and
nanoda configurations. PDF generation supplies presentation assurance only.

\begin{thebibliography}{99}
\bibitem{Guersenzvaig}
N. Guersenzvaig, \emph{Elementary criteria for irreducibility of $f(X^n)$},
2013, \url{https://arxiv.org/abs/1303.5333}.
\bibitem{HJ}
J. Harrington and L. Jones, \emph{Monogenic binomial compositions},
Taiwanese J. Math. 24 (2020), 1073--1090.
\url{https://doi.org/10.11650/tjm/200201}.
\bibitem{JKY}
A. Jakhar, R. Kalwaniya and P. Yadav,
\emph{A study of monogenity of binomial composition},
Acta Arith. 221 (2025), 387--395;
preprint \url{https://arxiv.org/abs/2402.10131}.
\url{https://doi.org/10.4064/aa250406-17-6}.
\bibitem{CPY}
A. Choudhary, S. Pisolkar and P. Yadav,
\emph{Power integral bases in polynomial compositions}, 2026,
\url{https://arxiv.org/abs/2606.09739}.
\bibitem{Gaal}
I. Ga\'al, \emph{On the monogenity of certain binomial compositions},
JP J. Algebra Number Theory Appl. 57 (2022), 1--16.
\url{https://doi.org/10.17654/0972555522026}.
\bibitem{Banerjee}
P. Banerjee, \emph{Variations on a theorem of Capelli},
Int. J. Number Theory 20 (2024), 893--913.
\url{https://doi.org/10.1142/S1793042124500465}.
\bibitem{Butler}
M. C. R. Butler, \emph{Reducibility criteria for polynomials of two general classes},
Proc. London Math. Soc. s3-7 (1957), 63--74.
\url{https://doi.org/10.1112/plms/s3-7.1.63}.
The complete article remains an explicit full-text access gap in the
present literature review.
\end{thebibliography}
\end{document}
